\section{特征值与特征向量的微分}
\label{sec:multivariable-eigen-differential}

特征值和特征向量由方程$\bm S\bm v=\lambda\bm v$确定，并不是矩阵元素的显式函数，前一节的显式计算不能直接使用。第~\ref{chap:condition-and-spectral-perturbation}章已经给出有限扰动下的稳定性界：对称扰动$\bm E$使每个排序后的特征值至多移动$\lVert\bm E\rVert_2$（Weyl界，定理~\ref{thm:ls-weyl-eigenvalue}），特征子空间的偏转受$\lVert\bm E\rVert_2$与谱隙之比控制（第~\ref{sec:ls-subspace-sensitivity}节）。这些界对一切方向取最坏情形。

本节回答另一个问题：沿一个指定的扰动方向，特征值与特征向量以什么速率变化，这个变化率在多大范围内可以用来预测。先用第~\ref{chap:implicit-inverse-functions}章的隐函数定理证明简单特征值所在的分支可微，再分别计算特征值与特征向量的微分，最后说明一阶近似的适用条件。

\subsection{特征值与特征向量的可微性}
\label{sec:multivariable-eigen-differentiability}

设实对称矩阵$\bm S_0\in\mathbb R^{d\times d}$有一个简单特征值$\lambda_0$，对应单位特征向量$\bm v_0$。简单是指这个特征值的代数重数为$1$；对实对称矩阵，这等价于其特征空间是一维的。即使规定长度为$1$，$\bm v_0$和$-\bm v_0$仍代表同一方向。为了在局部选定一支，我们要求附近的向量$\bm v$满足$\bm v_0^\top\bm v>0$。这里选出的只是一个局部分支，单位长度本身并不保证全局唯一。

把对称矩阵$\bm S$作为参数，选取其上三角元素作为独立坐标，便可把对称矩阵空间视为有限维欧氏空间。把向量$\bm v$和标量$\lambda$作为未知量，特征方程与单位长度条件合起来写成
\[
 F(\bm S,\bm v,\lambda)=
 \begin{pmatrix}\bm S\bm v-\lambda\bm v\\ (\bm v^\top\bm v-1)/2\end{pmatrix}=\bm0.
\]
这是$d+1$个方程、$d+1$个未知量。在$(\bm S_0,\bm v_0,\lambda_0)$处，关于$(\bm v,\lambda)$的导数作用于$(\bm h,\eta)$，其矩阵为
\[
 \begin{pmatrix}\bm S_0-\lambda_0\bm I&-\bm v_0\\
 \bm v_0^\top&0\end{pmatrix}.
\]
要用隐函数定理，需验证这个方阵可逆，即它只把零向量映为零。设它把$(\bm h,\eta)$映为零，即
\[
 (\bm S_0-\lambda_0\bm I)\bm h-\eta\bm v_0=\bm0,\qquad \bm v_0^\top\bm h=0 .
\]
先用$\bm v_0^\top$左乘第一式，利用对称性及$\bm S_0\bm v_0=\lambda_0\bm v_0$，左端第一项为零，得到$\eta=0$。于是$(\bm S_0-\lambda_0\bm I)\bm h=\bm0$，再由第二式，$\bm h$既属于$\lambda_0$的特征空间，又与$\bm v_0$正交。由于特征值简单，只能有$\bm h=\bm0$。

因此这个方阵可逆，隐函数定理保证在$\bm S_0$附近存在可微的特征值和单位特征向量分支；上述内积为正的条件固定了局部符号。

这里固定符号的条件不能省略。若不选定符号，同一个特征方向可以由$\bm v_0$或$-\bm v_0$表示，沿扰动路径任意切换符号会让“特征向量”看起来不连续；这种跳跃来自表示的不唯一，而不是特征方向本身的变化。选定$\bm v_0^\top\bm v>0$以后，隐函数定理给出的分支在邻域内唯一，可微性才有确定的含义。

\subsection{特征值的微分}
\label{sec:multivariable-eigenvalue-differential}

有了可微分支，就可以求它对每个小的对称扰动方向的一阶响应。设对称扰动方向$\bm E$使特征值与特征向量分别产生一阶变化$\mathrm d\lambda(\bm E)$和$\mathrm d\bm v(\bm E)$。对特征方程与归一化条件求微分，并在基点取值，得到
\begin{equation}
 \begin{gathered}
(\bm S_0-\lambda_0\bm I)\,\mathrm d\bm v(\bm E)
 =-\bigl(\bm E-\mathrm d\lambda(\bm E)\bm I\bigr)\bm v_0,\\
\bm v_0^\top\mathrm d\bm v(\bm E)=0.
 \end{gathered}
 \label{eq:multivariable-eigen-differential-system}
\end{equation}
用$\bm v_0^\top$左乘第一式，左侧为零，于是
\begin{equation}
 \mathrm d\lambda(\bm E)=\bm v_0^\top\bm E\bm v_0.
\label{eq:multivariable-eigenvalue-differential}
\end{equation}
这里的微分作用于全部小的对称方向；可微分支的存在由上一小节的隐函数定理保证，对特征方程形式求导本身不能承担这一步。

公式说明，特征值的变化率只看扰动在该特征方向上的读数$\bm v_0^\top\bm E\bm v_0$。由$\lVert\bm v_0\rVert_2=1$与谱范数的定义，
\[
 |\mathrm d\lambda(\bm E)|=|\bm v_0^\top\bm E\bm v_0|\leq\lVert\bm E\rVert_2 ,
\]
等号在$\bm E=\bm v_0\bm v_0^\top$时取到，所以简单特征值分支的局部绝对条件数（谱范数下）为$1$。与扰动方向正交的部分，例如满足$\bm E\bm v_0=\bm0$的扰动，一阶上不改变这个特征值。
\begin{example}[特征值沿两个方向的变化]
\label{ex:multivariable-eigenvalue-direction-check}
取$\bm S_0=\operatorname{diag}(2,1)$，简单特征值$\lambda_0=2$，$\bm v_0=\bm e_1$。先沿$\bm E_1=\bm e_1\bm e_1^\top$扰动，此时$\bm S_0+t\bm E_1=\operatorname{diag}(2+t,1)$，特征值为$2+t$，变化率为
\[
 1=\bm e_1^\top\bm E_1\bm e_1 ,
\]
余项为零。

再沿下面的方向扰动：
\[
 \bm E_2=\begin{bmatrix}0&1\\1&0\end{bmatrix}.
\]
公式给出$\mathrm d\lambda(\bm E_2)=\bm e_1^\top\bm E_2\bm e_1=0$。直接计算，$\bm S_0+t\bm E_2$的较大特征值为
\[
 \lambda(t)=\frac{3+\sqrt{1+4t^2}}2=2+t^2+O(t^4),
\]
变化确实是二阶的。

两个方向的谱范数都是$1$，Weyl界对两者都只给出$|\lambda(t)-2|\leq|t|$。
\end{example}

例子显示了两类结论的分工。Weyl界是对一切方向、任意步长都成立的最坏偏差；微分是指定方向上的变化率，它可以远小于最坏情形，甚至为零，但只在步长足够小时描述实际变化。用变化率预测有限扰动，还要知道余项多大，这正是最后一小节讨论的问题。

\subsection{特征向量的微分}
\label{sec:multivariable-eigenvector-differential}

特征向量的一阶变化由同一组方程~\eqref{eq:multivariable-eigen-differential-system}确定。把式~\eqref{eq:multivariable-eigenvalue-differential}给出的$\mathrm d\lambda(\bm E)$代回其中的第一式后，第一式右侧与$\bm v_0$正交。由于$\lambda_0$简单，$\bm S_0-\lambda_0\bm I$限制在$\bm v_0$的正交补上可逆，因此第二式选出了唯一的$\mathrm d\bm v(\bm E)$。

\paragraph{谱展开与谱隙}
在$\bm S_0$的标准正交特征基中，向量响应可以显式写出。设$\bm S_0\bm v_j=\lambda_j\bm v_j$（$j=1,\ldots,d$），$\bm v_1,\ldots,\bm v_d$标准正交，并记$\bm v_0=\bm v_i$、$\lambda_0=\lambda_i$。$\mathrm d\bm v(\bm E)$与$\bm v_i$正交，可写成$\sum_{j\neq i}c_j\bm v_j$。代入第一式，左侧为
\[
 (\bm S_0-\lambda_i\bm I)\sum_{j\neq i}c_j\bm v_j=\sum_{j\neq i}c_j(\lambda_j-\lambda_i)\bm v_j ;
\]
与$\bm v_j$（$j\neq i$）作内积，右侧给出$-\bm v_j^\top\bm E\bm v_i$，所以
\begin{equation}
 \mathrm d\bm v(\bm E)=\sum_{j\neq i}\frac{\bm v_j^\top\bm E\bm v_i}{\lambda_i-\lambda_j}\,\bm v_j .
\label{eq:multivariable-eigenvector-differential}
\end{equation}
设$d\geq2$，令谱隙$\gamma=\min_{j\neq i}|\lambda_i-\lambda_j|$，它因特征值简单而为正。由标准正交展开的Parseval等式，
\[
 \lVert\mathrm d\bm v(\bm E)\rVert_2^2=\sum_{j\neq i}\frac{(\bm v_j^\top\bm E\bm v_i)^2}{(\lambda_i-\lambda_j)^2}
 \leq\frac1{\gamma^2}\sum_{j}(\bm v_j^\top\bm E\bm v_i)^2=\frac{\lVert\bm E\bm v_i\rVert_2^2}{\gamma^2}\leq\frac{\lVert\bm E\rVert_2^2}{\gamma^2},
\]
即$\lVert\mathrm d\bm v(\bm E)\rVert_2\leq\lVert\bm E\rVert_2/\gamma$。$d=1$时单位特征向量分支恒为$\bm v_0$，微分为零。

特征值的导数界$|\mathrm d\lambda(\bm E)|\leq\lVert\bm E\rVert_2$与这里的特征向量导数界可以一起与第~\ref{sec:ls-spectral-value-sensitivity}节和第~\ref{sec:ls-subspace-sensitivity}节比较。Weyl界与谱隙界控制的是有限扰动下特征值和特征子空间的整体变化；这里给出的是带方向的一阶响应。特征值的响应$\bm v_i^\top\bm E\bm v_i$只看扰动在该特征方向上的读数；特征向量的响应以$\lambda_i-\lambda_j$为分母，说明谱隙小时，扰动可以把特征向量大幅地转向相邻的特征方向。
\begin{example}[特征向量的转向与谱隙]
\label{ex:multivariable-eigenvector-gap-check}
延续示例~\ref{ex:multivariable-eigenvalue-direction-check}，沿$\bm E_2$，公式给出
\[
 \mathrm d\bm v(\bm E_2)=\frac{\bm e_2^\top\bm E_2\bm e_1}{2-1}\bm e_2=\bm e_2 ,
\]
即$\bm v(t)\approx\bm e_1+t\bm e_2$。直接计算，较大特征值的特征向量与$(1,(\lambda(t)-2)/t)^\top$同向，而$(\lambda(t)-2)/t=t+O(t^3)$，与一阶预测一致。

把基点换成$\bm S_0=\operatorname{diag}(1+\gamma,1)$（$\gamma>0$），同一方向的响应变为$\mathrm d\bm v(\bm E_2)=\bm e_2/\gamma$：谱隙缩小为原来的几分之一，特征向量的转向速率就放大几倍。
\end{example}

微分给出的是指定方向上的转向速率，有限扰动下子空间之间的角度、投影距离和偏转界仍以第~\ref{chap:condition-and-spectral-perturbation}章为准：定理~\ref{thm:ls-davis-kahan}给出$\sin\theta\leq\lVert\bm E\rVert_2/\gamma$型的界，其中谱隙按该定理的规定在原矩阵与扰动后矩阵的相应特征值之间度量。导数界$\lVert\bm E\rVert_2/\gamma$与这个有限界形式相同，说明谱隙同时控制无穷小转向和有限偏转。

\subsection{一阶近似的适用条件}
\label{sec:multivariable-eigen-first-order-validity}

前两小节的公式只在简单特征值附近成立，而且只描述步长趋零时的行为。用它们预测有限扰动，需要扰动相对谱隙足够小。示例~\ref{ex:multivariable-eigenvector-gap-check}的第二个基点可以精确算出这一点。

这时扰动后的矩阵为
\[
 \bm S_0+t\bm E_2=\begin{bmatrix}1+\gamma&t\\t&1\end{bmatrix}.
\]
它的较大特征值的特征向量与$\bm e_1$的夹角$\theta$满足
\[
 \tan2\theta=\frac{2t}\gamma ,
\]
一阶预测则是$\theta\approx t/\gamma$。二者在$|t|\ll\gamma$时一致。

$|t|$与$\gamma$同阶时，实际角度明显小于预测，例如$t=\gamma/2$时$\theta=\pi/8\approx0.39$，预测为$0.5$。固定$t$而令$\gamma\to0$时，实际角度趋于$\pi/4$，预测却趋于无穷。

特征值也有同样的现象。较大特征值及其在$|t|\ll\gamma$时的近似为
\[
 1+\frac\gamma2+\sqrt{\frac{\gamma^2}4+t^2}\approx1+\gamma+\frac{t^2}\gamma ;
\]
一阶项为零，二阶项的系数$1/\gamma$随谱隙缩小而增大。所以一阶预测可靠的条件不是“扰动小”，而是“扰动相对谱隙小”。

上面的结论依赖所选特征值简单。重特征值使关于未知量的导数不可逆，因此不能直接沿用这个构造；但这只说明这条论证失效，并不断言每一条特征值分支都不可微，也不意味着矩阵多项式或可逆点处的求逆映射失去光滑性。

第~\ref{sec:multivariable-local-condition-number}节的$\operatorname{diag}(t,-t)$说明，在重特征值处，排序后的特征值可以有有限条件数而不可微，未排序的分支却仍然光滑。研究重特征值附近时，应改为考察相应的特征子空间，而不是强行指定其中某一个特征向量；谱簇子空间的稳定性已在第~\ref{sec:ls-subspace-sensitivity}节处理，这里不重复谱扰动界。

总之，本节只界定简单谱分支何时可微，以及微分公式何时可以用作一阶近似。重根附近的稳定性需要把相邻特征值合成隔离的特征簇，比较整个特征子空间；这部分直接使用第~\ref{chap:condition-and-spectral-perturbation}章的结论。
